% Selected general mathematical lemmas. Original line ranges are marked below.
% This is a source companion, not a standalone LaTeX document.

% Original source lines 269--424
\section{The four-extreme-element covering count}\label{sec:cover}

For a finite set $F$, put
\[
 R(F)=\{x+y:x,y\in F,\ x\ne y\},
 \qquad \Delta(F)=\card{2F\setminus R(F)}.
\]
Every value in $2F\setminus R(F)$ is a diagonal value $2x$, and
$x\mapsto2x$ is injective.  Hence
\begin{equation}
 \Delta(F)\leq\card F.
 \label{eq:delta-trivial}
\end{equation}

From a feasible set $A$ containing zero, take the four largest positive
elements and the four smallest negative elements, taking every element on a
side that has fewer than four.  Call their union $B$, and put
\[
 D=A\setminus B,\qquad d=\card D.
\]
Then $0\in D$ and $\card B\leq8$.

\begin{lemma}[Four-extreme-element covering]\label{lem:cover-count}
If $\FS(A)\subseteq4A$, then
\begin{equation}
 \boxed{\card{\FS(D)}\leq2\card{2D}+18d+85.}
 \label{eq:sharp-cover}
\end{equation}
\end{lemma}

\begin{proof}
We first work on the positive side.  Let
$p_1>p_2>p_3>p_4>0$ be the four largest positive elements, and put
$T_p=p_1+p_2+p_3+p_4$.  If $x\in\FS(D)$ and $x>0$, then
$T_p+x\in4A$.  Write a representation in decreasing order:
\[
 T_p+x=z_1+z_2+z_3+z_4,
 \qquad z_1\geq z_2\geq z_3\geq z_4.
\]
Classifying by the first failure of $z_i\leq p_i$ gives only
\[
\begin{array}{c|l}
2&2p_1+u+v,\\
3&p_1+2p_2+u\quad\text{or}\quad3p_2+u,\\
4&p_1+p_2+2p_3,\ p_1+3p_3,\ 2p_2+2p_3,\ p_2+3p_3,\ 4p_3.
\end{array}
\]
Subtracting $T_p$ and substituting $A=D\sqcup B$ gives the table
\[
\begin{array}{c|ccc}
\toprule
\text{source}&\text{translates of }2D&\text{translates of }D&\text{fixed slots}\\
\midrule
(2p_1-T_p)+A+A&1&8&36\\
(p_1+2p_2-T_p)+(A\setminus\{p_1\})&0&1&7\\
(3p_2-T_p)+(A\setminus\{p_1\})&0&1&7\\
\text{five fixed shapes}&0&0&5\\
\midrule
\text{total}&1&10&55\\
\bottomrule
\end{array}
\]
The negative side gives the same table.  Without deletions, the two sides and
zero yield
\begin{equation}
 \card{\FS(D)}\leq2\card{2D}+20d+111.
 \label{eq:raw-cover}
\end{equation}

If one side has fewer than four elements, all elements of that sign lie in
$B$, and $D$ has no subset sum of that sign.  The single remaining table is
already stronger than \eqref{eq:sharp-cover}.  We may therefore assume that
both sides have four extreme elements.  Write the absolute values on the
negative side as $q_1>q_2>q_3>q_4>0$.

Among the twenty translates of $D$, Appendix~\ref{app:slots} gives two
distinct pairs of centers $u\leq v$.  For any such pair,
\begin{equation}
 \card{(u+D)\cap(0,\infty)}+
 \card{(v+D)\cap(-\infty,0)}\leq d,
 \label{eq:halfline}
\end{equation}
because for one $z\in D$ the conditions $z>-u$ and $z<-v$ cannot both hold.
The two pairs reduce the linear contribution from $20d$ to $18d$.

There are $110$ fixed slots in the two tables.  Appendix~\ref{app:slots}
places them in one common universe and gives $26$ mutually disjoint groups.
The total score in every group is coordinatewise nonpositive in eight
strictly positive gap variables, so at least one slot in each group lies on
the wrong open half-line.  At most $110-26=84$ fixed slots survive.  Adding
zero gives the constant $85$, proving \eqref{eq:sharp-cover}.
\end{proof}

\section{Two four-element peeling steps}\label{sec:peeling}

\begin{lemma}[Four-element peeling]\label{lem:peeling}
Let $F$ be a finite real set containing zero and at least four nonzero
elements.  Four nonzero elements can be deleted to obtain $F'$ such that
\[
 \card{\FS(F)}\geq\card{\FS(F')}+\card{R(F)}-1.
\]
\end{lemma}

\begin{proof}
Choose the two largest positive and the two smallest negative elements,
filling with as-yet unselected nonzero elements if one side has fewer than
two.  Call the selected set $H$, put $F'=F\setminus H$, and define
\[
 \alpha=\sum_{h\in H,\,h<0}h,
 \qquad \beta=\sum_{h\in H,\,h>0}h,
 \qquad t=\beta-\alpha>0.
\]
The extreme choice gives $R(F)\subseteq[\alpha,\beta]$, while
\[
 \alpha+\FS(F')\subseteq\FS(F),
 \qquad \beta+\FS(F')\subseteq\FS(F).
\]
In every residue class modulo $t$ occupied by $\FS(F')$, the union of these
two translates contains at least one more point than $\FS(F')$.  In the
length-$t$ interval $[\alpha,\beta]$, each residue class occurs at most once,
except that the endpoints have the same residue.  Thus $R(F)$ occupies at
least $\card{R(F)}-1$ residue classes.  Add one point of $R(F)$ from every
class not already supplied by the two translates.  The stated inequality
follows.
\end{proof}

If $d\leq8$, then $\card A\leq d+8\leq16$, so the desired upper bound is
already trivial.  Hence assume $d\geq9$.  The set $D$ then has at least eight
nonzero elements, and after the first deletion the remaining set still has at
least four.  We may therefore apply the lemma twice:
\[
 D\supset D_1\supset E,
 \qquad \card{D_1}=d-4,
 \qquad \card E=d-8,
 \qquad 0\in E.
\]
Deleting four elements affects at most
\[
 4(d-4)+\binom42=4d-10
\]
unordered pairs of distinct elements.  Hence
\[
 \card{R(D_1)}\geq\card{R(D)}-(4d-10).
\]
If $s=\card{\FS(E)}$, the two peeling inequalities give
\begin{equation}
 \card{\FS(D)}\geq s+2\card{R(D)}-4d+8.
 \label{eq:peel-lower}
\end{equation}
Since $\card{2D}=\card{R(D)}+\Delta(D)$, comparison with
\eqref{eq:sharp-cover} eliminates $R(D)$ and gives
\begin{equation}
 \boxed{s\leq22d+2\Delta(D)+77.}
 \label{eq:core-upper}
\end{equation}


% Original source lines 438--626
\section{Subset-sum stability, increments, and defects}\label{sec:upper}
For an $r$-element positive real set $X$, define its subset-sum excess and its
doubled-value defect by
\[
 \mu(X)=\card{\FS(X)}-\left(1+\frac{r(r+1)}2\right),
 \qquad
 \delta_0(X)=\card{\{x\in X:2x\notin R(X\cup\{0\})\}}.
\]
The elementary increasing-chain argument gives $\mu(X)\geq0$.

We use the following published stability theorem in precisely its stated
range: if $r\geq4$, $0\leq M\leq r-4$, and
\[
 \card{\FS(X)}\leq1+\frac{r(r+1)}2+M,
\]
then $X=cB$ for some $c>0$ and an $r$-element set
$B\subset\mathbb Z_{>0}$ satisfying
\[
 \sum_{b\in B}b\leq\frac{r(r+1)}2+M.
\]
This is Theorem~1.1 of Carpenter--Defant--Kravitz~\cite{CDK2026}.

We also use Lemma~3.6 and Theorem~1.5 of
Balandraud--Girard--Griffiths--Hamidoune~\cite{BGGH2013}.  For a finite
subset of a finite abelian group, these give
\begin{equation}
 U\cap(-U)=\varnothing\quad\Longrightarrow\quad
 \card{\FS(U)}\geq2\card U,
 \label{eq:antisymmetric-bound}
\end{equation}
and, if $S=-S$, $\card S=2p\geq4$, and $\FS(S)$ is aperiodic, then
\begin{equation}
 \card{\FS(S)}\geq p(p-1)+5.
 \label{eq:symmetric-bound}
\end{equation}
Here a finite set is \emph{aperiodic} if only translation by zero fixes it.

\begin{lemma}[Aperiodic subset sums]\label{lem:aperiodic-linear}
Let $G$ be an abelian group and let
$R\subseteq G\setminus\{0\}$ have $m\geq2$ elements.  If $\FS(R)$ is
aperiodic, then
\begin{equation}
 \boxed{\card{\FS(R)}\geq2m-2.}
 \label{eq:aperiodic-linear}
\end{equation}
\end{lemma}

\begin{proof}
First suppose that $G$ is finite.  No member of $R$ has order two: otherwise
$\FS$ of that singleton, and hence $\FS(R)$, has a nonzero period.  Decompose
\[
 R=S\sqcup U,\qquad S=-S,\quad \card S=2p,\quad
 U\cap(-U)=\varnothing,\quad \card U=u.
\]
The subset-sum set of every subset of $R$ is aperiodic, since any of its
periods would also be a period after addition of the complementary
subset-sum set.  Kneser's theorem, whose stabilizer is therefore trivial,
gives
\[
 \card{\FS(R)}\geq\card{\FS(S)}+\card{\FS(U)}-1.
\]
If $p=0$, \eqref{eq:antisymmetric-bound} gives the stronger bound $2m$.
If $p=1$, then $\card{\FS(S)}=3$, and the right side is at least
$2u+2=2m-2$ (with the empty-$U$ case even larger).  If $p\geq2$ and
$u\geq1$, equations \eqref{eq:antisymmetric-bound} and
\eqref{eq:symmetric-bound} give
\[
 p(p-1)+5+2u-1
 =2m-2+(p-2)(p-3)\geq2m-2.
\]
If $p\geq2$ and $u=0$, the difference between the bound
$p(p-1)+5$ and $2m-2=4p-2$ is $p^2-5p+7>0$.

For a general $G$, work in the finitely generated subgroup generated by
$R$ and put $K=\FS(R)$.  Residual finiteness supplies a finite quotient whose
map is injective on $K+K-K$.  It also preserves aperiodicity: a nonzero period
of the image of $K$ can be written as the image of $k_1-k_0$; injectivity on
$K+K-K$ then lifts every translated point back to $K$, producing a nonzero
period of $K$.  The finite-group result in this quotient proves
\eqref{eq:aperiodic-linear}.
\end{proof}

Order a positive set as $X_i=\{x_1<\cdots<x_i\}$ and put
\[
 d_i=\card{(x_i+\FS(X_{i-1}))\setminus\FS(X_{i-1})},
 \qquad y_i=d_i-i.
\]
The standard chain argument gives $y_i\geq0$ and
\begin{equation}
 \mu(X_i)=\sum_{j=1}^i y_j.
 \label{eq:increment-sum}
\end{equation}
We shall also use the exact equality case in Lemma~2.1 of
Carpenter--Defant--Kravitz~\cite{CDK2026}: for $i\geq4$,
\begin{equation}
 y_i=0\quad\Longleftrightarrow\quad
 X_i=t\{1,2,\ldots,i\}\quad\text{for some }t>0.
 \label{eq:zero-increment}
\end{equation}

\begin{lemma}[Unit-increment inverse lemma]\label{lem:unit-increment}
If $i\geq6$ and $y_i=1$, then there are $t>0$ and $q\in\{1,\ldots,i\}$
such that
\begin{equation}
 X_i=t\bigl(\{1,2,\ldots,i+1\}\setminus\{q\}\bigr).
 \label{eq:unit-increment}
\end{equation}
\end{lemma}

\begin{proof}
Put $x=x_i$, reduce $x_1,\ldots,x_{i-1}$ modulo $x$ in the circle group
$\mathbb R/x\mathbb Z$, and call the resulting $(i-1)$-element set $R$.
Let $\rho=\card{\FS(R)}$.  Distinct occupied residues give distinct maximal
$x$-chains in $\FS(X_{i-1})$, while
$0,x_1,\ldots,x_{i-1}$ occupy distinct residues.  Hence
\[
 i\leq\rho\leq d_i=i+1.
\]
If $\FS(R)$ were aperiodic, Lemma~\ref{lem:aperiodic-linear} would give
$\rho\geq2i-4>i+1$, a contradiction.

Let $H$ be the nontrivial stabilizer of $\FS(R)$, write
$h=\card H$ and $\rho=hq_0$, and put $v=\card{R\cap H}$.  Kneser's theorem
applied to the $(i-1)$ sets $\{0,r\}$ gives
\[
 \rho\geq h(i-v),\qquad v\leq h-1,\qquad h+q_0\geq i+1.
\]
If $q_0\geq2$, then
\[
 \rho=hq_0\geq2(h+q_0)-4\geq2i-2>i+1,
\]
again impossible.  Thus $q_0=1$, so $\FS(R)=H$ and $R\subseteq H$.
The circle group has a unique subgroup of order $h$; with $t=x/h$ and
representatives in $(0,x)$,
\[
 X_i=t\bigl(\{1,\ldots,h\}\setminus D\bigr),\qquad
 \card D=h-i,\qquad i\leq h\leq i+1,
\]
where $D\subseteq\{1,\ldots,h-1\}$.  The case $h=i$ is the homogeneous
progression and has $y_i=0$.  Therefore $h=i+1$ and $D=\{q\}$, proving
\eqref{eq:unit-increment}.
\end{proof}

\begin{lemma}[Uniform defect versus excess]\label{lem:defect-excess}
For every finite positive real set $X$,
\begin{equation}
 \boxed{2\delta_0(X)-\mu(X)\leq10.}
 \label{eq:defect-excess}
\end{equation}
\end{lemma}

\begin{proof}
Use the prefixes above and write $\delta_i=\delta_0(X_i)$ and
$\Phi_i=2\delta_i-\mu(X_i)$.  Adding $x_i$ can remove old defects and can
create only the one new defect $x_i$, so
\begin{equation}
 \Phi_i\leq\Phi_{i-1}+2-y_i.
 \label{eq:defect-recurrence}
\end{equation}
For $i\leq5$, the trivial bounds $\delta_i\leq i$ and $\mu(X_i)\geq0$
give $\Phi_i\leq10$.

Proceed by induction for $i\geq6$.  If $y_i\geq2$, equation
\eqref{eq:defect-recurrence} preserves the bound.  If $y_i=0$, the published
increment-equality case \eqref{eq:zero-increment} makes $X_i$ a homogeneous
progression; then $\delta_i=1$, $\mu(X_i)=0$, and $\Phi_i=2$.

It remains to consider $y_i=1$.  By Lemma~\ref{lem:unit-increment}, after
scaling we have $X_i=[1,i+1]_{\mathbb Z}\setminus\{q\}$.  Every present
$k\in[2,i-1]_{\mathbb Z}$ has two disjoint representations
\[
 2k=(k-1)+(k+1)=(k-2)+(k+2),
\]
where $0$ is allowed; deleting one $q$ cannot destroy both.  The largest
element $i+1$ is defective.  The only possible additional defect is $1$
when $q=2$, or $i$ when $q=i-1$.  Thus $\delta_i\leq2$, and
$\delta_i=2$ is possible only for these two values of $q$.

For $q=2$ and $q=i-1$, the integer totals are respectively
\[
 T_i+i-1\quad\hbox{and}\quad T_i+2,
 \qquad T_i=\frac{i(i+1)}2.
\]
Both exceed $T_i+1$.  If $\mu(X_i)\leq1$, the published stability theorem
with $M=\mu(X_i)\leq1\leq i-4$ would give an integer normalization of total
at most $T_i+1$.  Since both displayed sets contain $1$, every integer
normalization has total at least the displayed total, a contradiction.
Hence $\mu(X_i)\geq2$ whenever $\delta_i=2$.  In all cases
$\Phi_i\leq2$, completing the induction and the proof.
\end{proof}

% Original source lines 679--934
\begin{lemma}[Small-increment inverse lemma]\label{lem:small-increment}
Let $i\geq6$ and $0\leq y_i\leq i-5$.  There are an integer $h$ and a set
$D_0\subseteq[1,h-1]_{\mathbb Z}$ such that, for some $t>0$,
\[
 i\leq h\leq i+y_i,\qquad \card{D_0}=h-i,\qquad
 X_i=t\bigl([1,h]_{\mathbb Z}\setminus D_0\bigr).
\]
\end{lemma}

\begin{proof}
Repeat the circle-group argument in Lemma~\ref{lem:unit-increment}, now with
$d_i=i+y_i$.  If the projected subset-sum set were aperiodic, its cardinality
$\rho$ would satisfy
\[
 \rho\geq2i-4>i+y_i.
\]
For a nontrivial stabilizer, retain the notation $h,q_0,v$ from that proof.
If $q_0\geq2$, then
\[
 \rho=hq_0\geq2(h+q_0)-4\geq2i-2>i+y_i.
\]
Thus $q_0=1$, the projected subset sums form the subgroup of order
$h=\rho$, and $i\leq h\leq i+y_i$.  Taking the representatives in $(0,x_i)$
gives the asserted interval with $h-i$ deletions.
\end{proof}

\section{Two stronger defect estimates}\label{sec:new-defect}
For a positive set $X$, write
\[
\Phi(X)=2\delta_0(X)-\mu(X),\qquad
\rhoD(X)=2-\Phi(X)=\mu(X)+2-2\delta_0(X).
\]
The uniform estimate already proved is $\Phi\leq10$.

\begin{lemma}[Deleted intervals]\label{lem:deleted}
Let $X=[1,r+q]_{\mathbb Z}\setminus D$, where $|D|=q$, $r+q\notin D$,
and $r\geq\max(5,4q+1)$. Then $\Phi(X)\leq2$.
For $q=1,2$ and $r\geq7$, one also has $\delta_0(X)\leq2q$.
\end{lemma}
\begin{proof}
The case $q=0$ is immediate. Put $h=r+q$ and
$C=\{h-d:d\in D\}$. The integer total excess is
\[
e=\sum X-T_r=\sum_{c\in C}c-T_{q-1}.
\]
There are $\min(k,h-k)$ disjoint representations of $2k$ by distinct
members of $[0,h]_{\mathbb Z}$ when $k\ne h/2$; for $k=h/2$, use $0+h$.
Therefore a defective present $k$ is either in $[1,q]$ or in $[r,h]$.
There are at most $q+1$ high defects. For $q\geq2$ there are at most $q-1$
low defects: having all of $1,\ldots,q$ present and defective would make
2 both absent (from defectivity of 1) and present. For $q=1$, the direct
one-deletion argument above gives at most two defects. These observations
prove $\delta_0(X)\leq2q$ and the last assertion.

Since $r>q$, no integer greater than one can divide all $r$ elements
of this interval with $q$ deletions; hence $X$ has gcd one.
If $\mu(X)\leq r-4$, the stability theorem in Section~\ref{sec:upper}
implies $\mu(X)=e$: any integer normalization of a primitive integer set
is an integer dilation, and the support lies in $[0,\sum X]$.
Suppose first that there are no low defects. Apart from $h$, let the $m$
high defects be $h-t$. Their presence and defectivity imply
$t\notin C$ and $2t\in C$. Replacing these distinct depths $2t$ by their
absent halves $t$ leaves $q$ distinct positive integers. Hence
\[
e\geq q+\sum t\geq q+m\geq2m=2\delta_0(X)-2.
\]
If there are $\ell\geq1$ low defects, their forced holes $2k$ have distinct
depths at least $r-q$, spaced by at least two. Thus
\[
e\geq\ell(r-q)+\ell(\ell-1)+T_{q-\ell}-T_{q-1}.
\]
Subtracting $2q+2\ell$ gives
\[
\ell(r-2q-2)+\tfrac32\ell(\ell-1)-q\geq0
\]
under $r\geq4q+1$. Since $\delta_0(X)\leq q+1+\ell$, this again implies
$e\geq2\delta_0(X)-2$.
Finally, if $\mu(X)>r-4$, then
\[
\Phi(X)\leq4q-(r-3)\leq2.
\]
All cases give the claim.
\end{proof}

\begin{theorem}[Threshold 14]\label{thm:phi14}
For every positive real set $X$ of size $r\geq14$,
\[
\boxed{\Phi(X)\leq2,\qquad\rhoD(X)\geq0.}
\]
\end{theorem}
\begin{proof}
For $9\leq i\leq14$, put $q_i=\lfloor(i-1)/4\rfloor$.
If $y_i\leq q_i$, the small-increment inverse lemma expresses $X_i$, up to
scaling, as an initial integer interval with at most $q_i$ deletions.
Lemma~\ref{lem:deleted} gives $\Phi(X_i)\leq2$.
If $y_i\geq q_i+1$, the defect recurrence gives
$\Phi(X_i)\leq\Phi(X_{i-1})+1-q_i$.
Starting from the uniform bound 10 at size 8 gives the following successive
upper bounds:
\[
\begin{array}{c|rrrrrrr}
i&8&9&10&11&12&13&14\\\hline
\Phi(X_i)\text{ upper bound}&10&9&8&7&6&4&2.
\end{array}
\]
After size 14, increments $y_i\geq2$ preserve the bound. For $y_i=0$ or 1,
the equality and unit-increment arguments in Section~\ref{sec:upper}
already give $\Phi(X_i)\leq2$.
\end{proof}

\begin{theorem}[Penalty for a non-progression support]\label{thm:nonap}
If $X$ is a positive real set of size $r\geq6$ and $\FS(X)$ is not an
arithmetic progression, then
\[
\boxed{\rhoD(X)\geq r-14.}
\]
\end{theorem}
\begin{proof}
At $r=6$ the uniform bound gives $\rhoD\geq-8=r-14$.
Induct on $r\geq7$. If $y_r\leq2$, the small-increment inverse lemma gives
a primitive deleted interval with at most two deletions. It has at most
four defects. Its non-progression support implies $\mu\geq r-3$ by the
stability theorem, since $\mu\leq r-4$ would force a full interval in the
primitive normalization. Consequently $\rhoD\geq r-9\geq r-14$.

Now let $y_r\geq3$. If the preceding support is not an arithmetic
progression, the induction hypothesis and
$\rhoD(X_r)\geq\rhoD(X_{r-1})+y_r-2$ give the result.
If the preceding support is $t[0,L]_{\mathbb Z}$, its translate by $x_r$
must be disjoint from it; an overlap would produce an arithmetic progression
again. Hence
\[
y_r=|\FS(X_{r-1})|-r\geq\frac{(r-1)(r-2)}2.
\]
Using $\rhoD(X_{r-1})\geq-8$ gives
$\rhoD(X_r)\geq (r-1)(r-2)/2-10\geq r-14$.
\end{proof}

\section{A sharper effective hole threshold}\label{sec:holes}
For a finite integer set $S$, its hull-hole count is
$\max S-\min S+1-|S|$. For a set on a specified lattice $u\mathbb Z$,
divide by $u$ before counting holes.

We use Lev's central-interval theorem~\cite{Lev1998}: if
$B\subseteq[1,L]_{\mathbb Z}$, $|B|=s$, and $L\leq3s/2-2$, then
\[
[2L-2s+1,\ \sum B-(2L-2s+1)]_{\mathbb Z}\subseteq\FS(B).
\]
Write $\mathcal S(L,H)=[0,L]_{\mathbb Z}\setminus(H\cup(L-H))$.
The endpoint bridge proved below states that an
integer prefix of size $m$ with support $\mathcal S(L,H)$ preserves the same
endpoint holes upon adjoining its next element with excess at most $E$, if
$2h<L$ and
\[
m\geq M(h,q,E):=
\max\left\{h,1+\left\lfloor
\frac{1+\sqrt{9+8(E+2q+2h)}}2\right\rfloor\right\},
\quad h=\max H,\quad q=|H|.
\]
Here $h=q=0$ is allowed. We need this explicit bridge, not an unspecified
``sufficiently large'' threshold from a classification theorem.

\begin{lemma}[Small increments with counted holes]\label{lem:counted-holes}
If $i\geq2E+5$ and $0\leq y_i=e\leq E$, then
\[
X_i=t([1,i+e]_{\mathbb Z}\setminus D),\qquad |D|=e,
\]
and, after division by $t$,
\[
\FS(t^{-1}X_i)=\mathcal S(L,H),\qquad
H\subseteq[1,2e]_{\mathbb Z},\qquad |H|\leq e.
\]
\end{lemma}
\begin{proof}
For $e=0$, the increment equality case gives the assertion directly.
Assume $e>0$. The small-increment inverse lemma first gives the same form with $\delta$
deletions, $0\leq\delta\leq e$, and largest coefficient $i+\delta$.
Remove that coefficient. The remaining $i-1$ integers lie in
$[1,i+\delta-1]$ and satisfy Lev's density condition, because
$i\geq2\delta+5$. Their support therefore has endpoint state
$H\subseteq[1,2\delta]$. Its total $L_0$ satisfies
\[
L_0\geq T_{i-1}>i+\delta+4\delta.
\]
Adding $i+\delta$ overlaps the two central intervals and leaves exactly
the same endpoint holes. Its support increment is $i+\delta$; comparison
with $i+e$ gives $\delta=e$.
Finally every lower hole in $[1,2e]$ is absent from the coefficient set
itself. Since $2e<i+e$ and there are exactly $e$ deleted coefficients,
$H\subseteq D$ and $|H|\leq e$. This last count halves the former bound
$|H|\leq2e$.
\end{proof}

\begin{theorem}[Effective hole bound]\label{thm:holes103}
Let $X$ be a positive real set of size $r\geq35$ with
$\FS(X)\subseteq u\mathbb Z$. If its hull-hole count $h$ satisfies
\[
14<h<|\FS(X)|-1,
\]
then
\[
\boxed{\mu(X)\geq r+103.}
\]
\end{theorem}
\begin{proof}
Suppose $\mu(X)\leq r+102$. Divide by $u$ to obtain an integer set $B$.
It is primitive: a gcd of at least two would leave at least one missing
lattice point between each two successive support values, contradicting
$h<|\FS(X)|-1$. The support is not a full integer interval.

We first obtain a tail of positive increments with a uniform budget.
If no prefix of size at least three has an arithmetic-progression support,
then $y_j\geq1$ for $j\geq4$, and
$\sum_{j=4}^r(y_j-1)\leq105$.
Otherwise let $k\geq3$ be the last size with such a support. The next
translate is disjoint, so
\[
y_{k+1}\geq T_k-k,
\qquad
\mu(B)\geq r-1+\frac{k(k-3)}2.
\]
Thus $k(k-3)/2\leq103$, whence $k\leq15$. All later increments are
positive, and
$\sum_{j=k+1}^r(y_j-1)\leq102+k\leq117$.
In either case there is $r_0\leq16$ such that
\begin{equation}
y_j\geq1\ (j\geq r_0),\qquad
\sum_{j=r_0}^r(y_j-1)\leq117.\label{eq:tail117}
\end{equation}


The window $19\leq j\leq r$ has at least 17 indices. If every increment
there were at least eight, it would contribute at least $7\cdot17=119$
to \eqref{eq:tail117}. Therefore some $i_0\geq19$ has $1\leq y_{i_0}=e\leq7$.
Lemma~\ref{lem:counted-holes} gives
\[
B_{i_0}=t([1,i_0+e]_{\mathbb Z}\setminus D),\quad
H\subseteq[1,14],\quad |H|\leq7.
\]
The coefficient set has gcd one since $i_0>e$, so $t$ is a positive integer.
If $t>1$, primitivity of $B$ gives a first later coefficient outside
$t\mathbb Z$. At that step the two supports are disjoint, and
\[
y_j-1\geq\frac{i_0(i_0-1)}2-1\geq170>117,
\]
contradicting \eqref{eq:tail117}. Hence $t=1$.

Every subsequent increment is at most 118, and
\[
M(14,7,118)=19,\qquad T_{19}=190>28.
\]
The endpoint bridge now preserves $H$ through the whole remaining tail.
The final number of holes is $2|H|\leq14$, contradicting $h>14$.
\end{proof}


% Original source lines 936--974
\begin{lemma}[Unpaired absolute values]\label{lem:unpaired}
Let $0\in F$, let $e=\card F$, and write the distinct nonzero absolute values
as $0<b_1<\cdots<b_r$.  Let $m_i\in\{1,2\}$ count how many of
$\{b_i,-b_i\}$ lie in $F$.  If
\[
 p=\card{\{i:m_i=2\}},\qquad q=\card{\{i:m_i=1\}},
\]
and
\[
 J=\card{\{(i,j):i<j,\ m_i=1,\ m_j=2\}},
\]
then
\begin{equation}
 \card{\FS(F)}\geq\frac{e^2+3+q^2}{4}+J.
 \label{eq:unpaired}
\end{equation}
\end{lemma}

\begin{proof}
Translate $\FS(F)$ by the sum of the absolute values of the negative members
of $F$.  The resulting set is
\[
 \left\{\sum_{i=1}^r k_i b_i:0\leq k_i\leq m_i\right\}.
\]
Before $b_i$ is introduced, the residues represented modulo $b_i$ by
$0,b_1,\ldots,b_{i-1}$ are distinct.  Extending each corresponding maximal
$b_i$-chain by the $m_i$ new coefficient levels creates at least $i m_i$
new sums.  Therefore
\[
 \card{\FS(F)}\geq1+\sum_{i=1}^r i m_i.
\]
If the paired positions are $s_1<\cdots<s_p$, then
\[
 \sum_{i=1}^r i m_i
 =\frac{r(r+1)}2+\sum_{j=1}^p s_j
 =\frac{r(r+1)}2+\frac{p(p+1)}2+J.
\]
Using $r=p+q$ and $e=2p+q+1$ gives \eqref{eq:unpaired}.
\end{proof}

% Original source lines 1241--1249
Stanchescu simultaneous inverse theorem: if finite integer sets $U,V$, each of
size at least two, satisfy
\[
 |U+V|<|U|+|V|+\min(|U|,|V|)-3,
\]
then, on their common difference lattice, their hull lengths are at most
$|U+V|-|V|+1$ and $|U+V|-|U|+1$. Stanchescu's corollary allows the slightly
larger threshold with $-2-\delta$, where $\delta$ is 1 for equal diameters
and 0 otherwise. Thus the strict hypothesis above is sufficient in both cases.

% Original source lines 1266--1321
\begin{lemma}[Elementary endpoint bridge]\label{lem:bridge-proof}
The endpoint bridge stated in Section~\ref{sec:holes} follows directly from
its displayed threshold $M(h,q,E)$.
\end{lemma}
\begin{proof}
Let $S=\mathcal S(L,H)$ be the support of $m$ positive integers. Then
$L\geq T_m$, $|S|=L+1-2q$, and $x>\max B\geq m\geq h$.
If $x\geq L-2h$, the part of $x+S$ above $L$ has at least
$\min(x,L+1)-2q\geq L-2h-2q$ points. But the definition of $M$ gives
\[
 T_m-2h-2q>m+1+E,
\]
contradicting the permitted increment. Thus $h<x<L-2h$.
The intervals $[h+1,L-h]$ and $[x+h+1,x+L-h]$ overlap and cover the
whole central interval $[h+1,L+x-h]$. Since $x>h$, the lower holes stay $H$
and the upper holes are their reflections about $(L+x)/2$.
Consequently $S\cup(x+S)=\mathcal S(L+x,H)$, its cardinality increases by
$x$, and $x=m+1+e_x$ follows. This also supplies a short self-contained proof
of the bridge used in the earlier hole estimate.
\end{proof}

\begin{lemma}[Parameterized effective hole test]\label{lem:hole-parameter}
Let $C\geq0$ and $b\geq1$ be integers. Define
\[
 K=\max\{k\geq3:k(k-3)/2\leq C+1\},\quad
 D=C+\max(3,K),\quad E=D+1.
\]
Suppose integers $J,r$ satisfy
\begin{gather*}
 J\geq\max\{K+1,2b+5,M(2b,b,E)\},\\
 (r-J+1)b>D,\qquad T_{J-1}-1>D.
\end{gather*}
If a positive $r$-element set has, on a specified lattice, a hull-hole count
$h$ satisfying $2b<h<|\FS|-1$, then $\mu\geq r+C+1$.
\end{lemma}
\begin{proof}
Assume instead $\mu\leq r+C$, and divide by the lattice step. The strict
upper bound on $h$ makes the integer set primitive, and its support is not
an arithmetic progression. The tail argument of
Theorem~\ref{thm:holes103}, with $102$ replaced by $C$, gives a starting
index at most $K+1$ (or 4 if no prefix of size at least three has progression
support) after which
\[
 y_j\geq1,\qquad \sum(y_j-1)\leq D.
\]
Here a last progression prefix of size $k$ costs at least
$r-1+k(k-3)/2$, so $k\leq K$; without such a prefix the budget is $C+3$.
The window $J,\ldots,r$ contains some increment $1\leq y_i\leq b$.
Lemma~\ref{lem:counted-holes} gives a deleted interval, with endpoint holes
in $[1,2b]$ and at most $b$ lower holes. Its scale is an integer. A first
later coefficient off this scale would cost
$y_j-1\geq T_{i-1}-1\geq T_{J-1}-1>D$; primitivity therefore makes the
scale one. Every later increment is at most $E$. The endpoint bridge
preserves the holes, since $i\geq J$ and $L\geq T_J>4b$.
The final hole count is at most $2b$, a contradiction.
\end{proof}

% Original source lines 3242--3278
\begin{lemma}\label{lem:six-flexible98}
Let $X$ be a positive real set of size $r\geq15$, and let
$\rho=\rhoD(X)$. If $0\leq\rho<5(r-14)$, then
\[
 \mu(X)\leq\rho+22+2\lfloor\rho/5\rfloor
 \leq\tfrac75\rho+22.
\]
Consequently, a lower bound $\mu(X)\geq u$ implies
\[
 \rho\geq\min\left(5(r-14),
              \left\lceil\frac{5(u-22)}7\right\rceil\right).
\]
\end{lemma}
\begin{proof}
Use the ordered prefixes $X_j$ and the increments $y_j$ from the defect
recurrence. The counting part of Lemma~\ref{lem:deleted}, independently of
its stronger hypothesis for the estimate on $\Phi$, gives
$\delta_0([1,j+q]\setminus H)\leq2q$ for $q\geq2$ and $j>q$.
Indeed every defect belongs to $[1,q]$ or $[j,j+q]$, and there are at most
$q-1$ defects in the former band: if all of $1,\ldots,q$ were present
and defective, defectivity of 1 would require 2 to be absent.
There are at most $q+1$ in the latter band. For $q=1$ the one-deletion
argument gives at most two defects (at the prefix sizes used here);
for $q=0$ there is one.

If every $y_j$ for $15\leq j\leq r$ were at least 7, nonnegativity of
$\rhoD(X_{14})$ and the recurrence would give
$\rho\geq5(r-14)$. Choose the largest $j\geq15$ with $y_j\leq6$.
The small-increment inverse lemma applies, since $6\leq j-5$, and makes
$X_j$, up to scaling, an interval with at most six deletions.
The preceding count gives $\delta_0(X_j)\leq12$.
Each subsequent step adds at most one defect, while all its increments
are at least 7; hence $r-j\leq\lfloor\rho/5\rfloor$ and
$\delta_0(X)\leq12+\lfloor\rho/5\rfloor$.
Now use $\mu=\rho-2+2\delta_0$. The final assertion follows by separating
$\rho\geq5(r-14)$ from its complement.
\end{proof}

% Original source lines 3692--3726

\begin{lemma}[Flexible defect bound with an integer inverse]\label{lem:flex96}
Let $X$ be a positive real set of size $r\geq15$, and let
$2\leq\beta\leq10$ be an integer. If
$0\leq\rhoD(X)<(\beta-1)(r-14)$, then
\[
 \mu(X)\leq\rhoD(X)+4\beta-2+
                    2\left\lfloor\frac{\rhoD(X)}{\beta-1}\right\rfloor.
\]
For an integer lower bound $\mu(X)\geq u$, put
$v=\max(0,u-4\beta+2)$ and write $v=(\beta+1)q+s$, $0\leq s\leq\beta$.
Then
\begin{equation}
 \rhoD(X)\geq\min\bigl((\beta-1)(r-14),
                  (\beta-1)q+\min(s,\beta-1)\bigr).
 \label{eq:flex96-inverse}
\end{equation}
\end{lemma}
\begin{proof}
If all increments $y_i$ for $15\leq i\leq r$ exceed $\beta$, the
recurrence $\rho_i\geq\rho_{i-1}+y_i-2$ contradicts the hypothesis.
Choose the last index $i\geq15$ with $y_i\leq\beta$.
Since $\beta\leq i-5$, Lemma~\ref{lem:small-increment} gives an interval
with at most $\beta$ deletions. The defect-counting argument in
Lemma~\ref{lem:six-flexible98} uses only $i>\beta$ and gives at most
$2\beta$ defects; the zero-deletion case has one defect.
Every subsequent step adds at most one defect and raises $\rho$ by at
least $\beta-1$. Nonnegativity at index $i$ therefore gives
$r-i\leq\lfloor\rhoD(X)/(\beta-1)\rfloor$. Use
$\mu=\rhoD-2+2\delta_0$.
Finally, the least nonnegative integer $t$ satisfying
$t+2\lfloor t/(\beta-1)\rfloor\geq v$ is
$(\beta-1)q+\min(s,\beta-1)$: each complete block of $\beta-1$ in $t$
adds $\beta+1$ to the left side. This proves the stated inverse, including
the alternative where the small-defect hypothesis fails.
\end{proof}

% Original source lines 4263--4376
\section{The fixed-slot deletion certificate}\label{app:slots}

Let
\[
 p_1>p_2>p_3>p_4>0,
 \qquad q_1>q_2>q_3>q_4>0,
\]
and define the eight positive gaps
\[
\begin{gathered}
 r_1=p_1-p_2,\ r_2=p_2-p_3,\ r_3=p_3-p_4,\ r_4=p_4,\\
 s_1=q_1-q_2,\ s_2=q_2-q_3,\ s_3=q_3-q_4,\ s_4=q_4.
\end{gathered}
\]
Put
\[
\begin{aligned}
 c_p&=p_1-p_2-p_3-p_4,&
 \alpha_p&=p_2-p_3-p_4,&
 \beta_p&=2p_2-p_1-p_3-p_4,\\
 c_q&=q_1-q_2-q_3-q_4,&
 \alpha_q&=q_2-q_3-q_4,&
 \beta_q&=2q_2-q_1-q_3-q_4.
\end{aligned}
\]
The two ordered eight-tuples are
\[
 B^+=(p_1,p_2,p_3,p_4,-q_1,-q_2,-q_3,-q_4),
\]
\[
 B^-=(q_1,q_2,q_3,q_4,-p_1,-p_2,-p_3,-p_4).
\]
For $\varepsilon\in\{+,-\}$, set
\[
 C^\varepsilon_{ij}=c_\varepsilon+B^\varepsilon_i+B^\varepsilon_j,
 \quad
 U^\varepsilon_i=\alpha_\varepsilon+B^\varepsilon_i,
 \quad
 V^\varepsilon_i=\beta_\varepsilon+B^\varepsilon_i,
\]
where $1\leq i\leq j\leq8$, and $U,V$ use $2\leq i\leq8$.  For a decreasing
quadruple $x=(x_1,x_2,x_3,x_4)$, define
\[
\begin{aligned}
 H_1(x)&=x_3-x_4,& H_2(x)&=2x_3-x_2-x_4,\\
 H_3(x)&=x_2+x_3-x_1-x_4,& H_4(x)&=2x_3-x_1-x_4,\\
 H_5(x)&=3x_3-x_1-x_2-x_4.
\end{aligned}
\]
Write $H_i^+=H_i(p)$ and $H_i^-=H_i(q)$.  The common slot universe has size
\[
 2\left(\binom{9}{2}+7+7+5\right)=110.
\]

The following table gives $15$ templates.  Every starred row is also used
after interchanging $+$ and $-$.  The four unstarred rows and the two
orientations of the eleven starred rows give $26$ groups.

\begingroup
\scriptsize
\setlength{\tabcolsep}{3.6pt}
\begin{longtable}{c>{\raggedright\arraybackslash}p{0.43\textwidth}p{0.39\textwidth}}
\toprule
&\textbf{sum of scores}&$(r_1,r_2,r_3,r_4;s_1,s_2,s_3,s_4)$\\
\midrule
\endfirsthead
\toprule
&\textbf{sum of scores}&$(r_1,r_2,r_3,r_4;s_1,s_2,s_3,s_4)$\\
\midrule
\endhead
1&$C^+_{25}+C^-_{25}$&$(0,0,-1,-2;0,0,-1,-2)$\\
$2^*$&$C^+_{48}+H^+_3$&$(0,0,0,-1;0,0,0,-1)$\\
3&$C^+_{58}+C^-_{58}$&$(0,-1,-2,-4;0,-1,-2,-4)$\\
$4^*$&$C^+_{77}+H^+_4$&$(0,-1,0,-2;0,0,-2,-2)$\\
5&$U^+_6+U^-_6$&$(0,0,-1,-2;0,0,-1,-2)$\\
6&$V^+_6+V^-_6$&$(-1,0,-1,-2;-1,0,-1,-2)$\\
$7^*$&$C^+_{15}+H^+_1+C^-_{55}$&$(0,-1,-1,-3;0,-1,-2,-3)$\\
$8^*$&$C^+_{26}+C^-_{57}+V^-_4$&$(0,0,-2,-3;0,0,-2,-3)$\\
$9^*$&$C^+_{35}+V^+_3+C^-_{46}$&$(0,0,0,-2;0,-1,-2,-2)$\\
$10^*$&$C^+_{36}+U^+_4+U^-_5$&$(0,0,-1,-2;0,0,-1,-2)$\\
$11^*$&$C^+_{44}+U^+_8+H^+_5$&$(0,-1,0,-1;0,0,0,-1)$\\
$12^*$&$C^+_{45}+V^+_8+C^-_{67}$&$(0,0,-3,-4;0,-1,-2,-4)$\\
$13^*$&$C^+_{47}+V^+_7+H^+_2$&$(0,0,0,-2;0,0,-2,-2)$\\
$14^*$&$C^+_{56}+V^+_2+C^-_{66}$&$(0,0,-2,-4;0,-2,-3,-4)$\\
$15^*$&$C^+_{68}+U^+_3+V^-_5$&$(0,0,-1,-3;-1,0,-1,-3)$\\
\bottomrule
\end{longtable}
\endgroup

Every gap is strictly positive, while each coefficient vector is
coordinatewise nonpositive and has a negative coordinate.  Hence all slots in
one group cannot have positive score.  The $26$ groups are mutually disjoint in the common universe, as can be checked directly from the displayed slot labels.

For the two translate pairings, first take
\[
 u_1=c_p-q_1,\qquad v_1=-c_q+p_1.
\]
Then
\[
 u_1-v_1=-(p_2+p_3+p_4+q_2+q_3+q_4)<0.
\]
If $\beta_p+\beta_q\leq0$, take
$u_2=\beta_p$ and $v_2=-\beta_q$.  Otherwise take
\[
 u_2=c_p-q_2,\qquad v_2=-c_q+p_2.
\]
The identity
\[
 u_2-v_2=-(\beta_p+\beta_q)-2(p_3+p_4+q_3+q_4)<0
\]
proves the needed order in the second branch.  The two pairs use distinct
table slots, so \eqref{eq:halfline} can be applied twice.

% Original source lines 5569--5598
\begin{thebibliography}{9}
\setlength{\itemsep}{0pt}
\bibitem{Stanchescu1996}
Y.~Stanchescu, \emph{On addition of two distinct sets of integers},
Acta Arithmetica \textbf{75} (1996), 191--194. In particular, Theorem~4
and Corollary~2 give the simultaneous hull bounds.
\url{https://doi.org/10.4064/aa-75-2-191-194}.

\bibitem{CDK2026}
R.~Carpenter, C.~Defant, and N.~Kravitz,
\emph{Sets with few subset sums}, arXiv:2605.05498 (2026).
\url{https://arxiv.org/abs/2605.05498}.

\bibitem{BGGH2013}
E.~Balandraud, B.~Girard, S.~Griffiths, and Y.~Ould Hamidoune,
\emph{Subset sums in abelian groups},
European Journal of Combinatorics \textbf{34} (2013), 1269--1286.
\url{https://arxiv.org/abs/1112.1929}.

\bibitem{LevSmeliansky1995}
V.~F.~Lev and P.~Y.~Smeliansky,
\emph{On addition of two distinct sets of integers},
Acta Arithmetica \textbf{70} (1995), 85--91.
\url{https://doi.org/10.4064/aa-70-1-85-91}.

\bibitem{Lev1998}
V.~F.~Lev, \emph{On consecutive subset sums},
Discrete Mathematics \textbf{187} (1998), 151--160.
\url{https://doi.org/10.1016/S0012-365X(98)80006-X}.

\end{thebibliography}
